% ============================================================
%  第七章
% ============================================================
\section{另一条线：几何 Langlands 与 Hecke 变换}
\markboth{第七章\ 另一条线：几何 Langlands 与 Hecke 变换}{第七章\ 另一条线：几何 Langlands 与 Hecke 变换}

\noindent{\small\color{shuyiGray}对应底本：原 p.37--p.40。}\par\medskip

\subsection{一条平行的故事线}

\srcp{37}在往下讲之前，我得先往回退一步，解释这个故事里的另一条线索：大约从 1990 年起，Alexander Beilinson 与 Vladimir Drinfeld 为普通的 Riemann 面（而不是数域）表述了 Langlands 对应的一个“几何”版本。从一开始，他们就大量使用了物理学里熟悉的构件，例如流代数（仿射 Kac--Moody 代数）。我做过许多努力去理解他们在说什么，但在那些年里，我觉得这不可能。看起来，物理学里熟悉的构件被以完全陌生的方式打乱重组了。

\subsection{几何 Hecke 变换}

\srcp{38}尤其让我困惑的，是几何 Hecke 变换（据我所知，最初是由 M.~S.~Narasimhan 与 S.~Ramanan 在 1970 年代引入的），而它在几何 Langlands 中是绝对核心的。设定如下。$E\to C$ 是 Riemann 面 $C$ 上的一个全纯 $G_\CC$-丛（$G_\CC$ 是单复 Lie 群）。$E$ 在点 $p\in C$ 处的一个\emph{Hecke 修改}，是一个全纯 $G_\CC$-丛 $E'\to C$，连同 $C\setminus p$ 上 $E$ 与 $E'$ 之间的一个同构
\[
  \varphi:\ E|_{C\setminus p}\ \xrightarrow{\ \cong\ }\ E'|_{C\setminus p}.
\]
例如，若 $G_\CC=\CC^*$，则 $E$ 是一个全纯线丛 $L\to C$；$L$ 在 $p$ 处的一个 Hecke 修改的例子是 $L'=L(p)=L\otimes\cO(p)$——按定义，它在离开 $p$ 的地方与 $L$ 同构。若 $G_\CC=SL(2,\CC)$，我们可以用秩 $2$ 的全纯向量丛表示 $E$；对 $E$ 的任一局部线丛分解 $L_1\oplus L_2$，$E$ 的一个 Hecke 修改的例子是 $E'=L_1(p)\oplus L_2(-p)$。

\begin{figure}[htbp]
\centering
\begin{tikzpicture}[scale=1.0]
  % 曲线 C
  \draw[thick] (-4.2,0) .. controls (-2,0.5) and (-1,-0.5) .. (0,0.1)
               .. controls (1,0.6) and (2,-0.4) .. (4.2,0.2);
  \node at (-4.6,0.1) {\small $C$};
  \fill (0,0.1) circle (1.8pt);
  \node[below=3pt] at (0,0.05) {\small $p$};
  % 两个丛
  \node at (-2.6,1.25) {\small $E$};
  \node at (-2.6,-1.15) {\small $E'$};
  \draw[->,thick] (-2.1,1.05) .. controls (-1.2,0.75) and (-1.2,0.55) .. (-0.25,0.32);
  \draw[->,thick] (-2.1,-0.95) .. controls (-1.2,-0.65) and (-1.2,-0.45) .. (-0.25,-0.22);
  % 同构箭头（避开 p）
  \draw[<->,dashed] (-0.55,0.55) -- (-0.55,-0.4);
  \node[right=2pt] at (-0.5,0.05) {\scriptsize $\varphi$};
  % p 处断开
  \draw[white,line width=3pt] (-0.12,-0.55) -- (0.12,0.75);
  \node[above] at (2.4,1.1) {\scriptsize 在 $C\setminus p$ 上同构};
\end{tikzpicture}
\caption{Hecke 修改的几何图像：两个丛只在 $p$ 点不同，在挖掉 $p$ 的曲线上则由 $\varphi$ 等同。所谓“修改”，是允许丛在一点处换一个“类型”。}
\label{fig:hecke}
\end{figure}

\subsection{类型与几何 Satake}

\srcp{39}在给定的点 $p$ 处，可能有不同类型的 Hecke 修改。例如，若 $G_\CC=\CC^*$，我们不必取 $L\to L(p)$，而可以取 $L\to L(np)=L\otimes\cO(p)^n$，$n$ 为任意整数。几何 Langlands 纲领中一项重要的输入，是\emph{几何 Satake 对应}（由 Lusztig、Drinfeld、Ginzburg、Mirković 与 Vilonen 等人发展）：它说，$G$ 的 Hecke 修改的可能类型，由对偶群 $\gvee$ 的不可约有限维表示来标号。（例如，若 $G_\CC=\CC^*$，对偶群还是 $\CC^*$；$\CC^*$ 的不可约表示由整数标号，这个整数就是 $L\to L(np)$ 里的 $n$。）

\subsection{与物理的接点}

\srcp{40}多年来，我一直琢磨：这样的几何 Hecke 变换究竟怎么可能与物理有关——因为我假定，如果几何 Langlands 要有物理解释，它们就非得有关不可。最终发现，用量子规范理论的基本构件可以给出一个简单的答案，那就是 Wilson 与 't Hooft 算子之间的电—磁对偶，以及 Bogomolny 方程——后者被 Atiyah、Hitchin 以及不少物理学家彻底研究过。（这个答案是我 2006 年与 Anton Kapustin 合写的、关于规范理论与几何 Langlands 的论文里的主要想法之一。今天没有时间解释其余的想法了。）

\begin{pizhu}{两套 Langlands 的词典}
把几何 Langlands 与数域上的经典 Langlands 并排看，会清楚得多。粗略的对照是：

\begin{center}\small
\begin{tabular}{@{}ll@{}}
\toprule
数域 $F$ & Riemann 面 $C$ \\
\midrule
Galois 群的表示 $\mathrm{Gal}(\bar F/F)\to\gvee$ & 局部系统 $\mathbb{L}$（$\gvee$-丛带平坦联络） \\
自守表示 $\pi$ & $\mathrm{Bun}_G$ 上的 Hecke 特征层 $\cF$ \\
Frobenius 的特征多项式 & Hecke 算子 $H_\lambda$ 的特征值 \\
\bottomrule
\end{tabular}
\end{center}

“特征层”的意思是：对每个支配余权 $\lambda$，Hecke 算子 $H_\lambda$ 作用在 $\cF$ 上等于乘以 $\mathbb{L}$ 的对应表示 $V_\lambda$。在数域情形，这就是“自守表示与 Galois 表示有相同的 $L$-函数”；在几何情形，它被提成了一条关于层与算子的精确陈述。
\end{pizhu}

\begin{pizhu}{Hecke 修改的“类型”}
在 $p$ 点附近取一个形式圆盘 $D$，把丛平凡化。一个 Hecke 修改在 $D$ 上就是给一个 $G_\CC$-值函数 $g(z)$（$z$ 是 $D$ 上的坐标），它在 $z\ne0$ 处有定义且在 $z=0$ 处可以有极点；两个这样的函数若在 $z\ne0$ 处相差一个全纯可逆的函数，就给出同一个修改。于是修改的“类型”由 $g$ 在 $z=0$ 处的\emph{估值}决定：在 $G$ 的极大环面 $T$ 上，$g(z)=z^{\lambda}$，其中 $\lambda$ 是余特征标格 $X_*(T)$ 里的一个元素，再模掉 Weyl 群的作用。因此
\[
  \{\text{Hecke 修改的类型}\}\ \cong\ X_*(T)/W\ =\ \{\text{支配余权}\}.
\]

对 $GL(n)$，$\lambda=(k_1\ge\dots\ge k_n)$ 是一组整数，修改就是把各个直和项的 twisting 数分别移动 $k_i$。对 $SL(2)$，$\lambda=n\ge0$，对应的正是 $\mathrm{Sym}^n$ 表示——这就是第 39 页里“$n$”的来源。
\end{pizhu}

\begin{pizhu}{几何 Satake 的精确陈述}
设 $G$ 是复约化群，$\cO=\CC[[z]]$、$\cK=\CC((z))$，仿射 Grassmannian 是 $\mathrm{Gr}_G=G(\cK)/G(\cO)$。它按 $G(\cO)$-轨道分解为 $\mathrm{Gr}_G=\bigsqcup_{\lambda\in X_*(T)/W}\mathrm{Gr}^\lambda$，每个轨道是有限维的（维数 $=\ip{2\rho}{\lambda}$）。

几何 Satake（Mirković--Vilonen，1999）断言：$\mathrm{Gr}_G$ 上 $G(\cO)$-等变反常层构成的范畴，与 $\gvee$ 的有限维表示范畴等价；轨道 $\mathrm{Gr}^\lambda$ 的 IC 层对应不可约表示 $V_\lambda$。

这就是“类型由对偶群的表示标号”的准确含义：Hecke 修改的几何类型（一个轨道）与对偶群的表示（一个 IC 层）是同一套组合数据的两个面貌。Witten 说的“打乱重组”，指的正是这件事——物理学家熟悉的仿射 Kac--Moody 表示论，在这里被重新组织成了层论。
\end{pizhu}
